\section{Bounds for the original replica metrics}
\label{sec:replica_pinned_upper_estimates}

The estimates below are the quadratic-form and inverse-power steps in
\cite[Section~8, lines~367--418]{OpenAI2026AreaLaw}. All replica weights
are computed from the original full system.

% Source: OpenAI/math revision adc7f1241b42e322a6451854ab7e4b4c146bf78a,
% 07-comparators.tex, comparator:pinned-metric, lines 299--308 and 405--415.
\begin{theorem}[Exponential upper bound for one pinned metric factor]
    \label{thm:replica_pinned_factor_exp_upper}
    \lean{TensorPower.exists_replicaPinnedFactor_le_exp}
    \leanok
    \uses{thm:replica_metric_basic,def:label_projectors}
    Fix a finite tensor product of nonzero finite-dimensional Hilbert
    spaces and $t>0$. On every subsystem $Q$, use the replica metric
    $W_{Q,k}$ defined by the common weights of the original full system,
    and put $F_{Q,k}=\sum_\lambda(\log d_\lambda)\pi_{Q,k}^\lambda$.
    There is $C\ge0$, depending on the fixed local spaces and $t$, such
    that for every integer $k\ge0$ and every pair of disjoint regions
    $U,Y$,
    \begin{align*}
        \bigl(W_{U,k}^{-1}W_{U\cup Y,k}^{-1}W_{Y,k}\bigr)^2
         & \le (k+2)^C
        \exp\bigl(2t(F_{U,k}+F_{U\cup Y,k}-F_{Y,k})\bigr).
    \end{align*}
    The inequality holds on the whole copy space. Empty regions and
    zero copies are included. With $a=2t$, it gives the exponential
    comparison in the proof of the sharp upper bound
    in \cite[Section~8]{OpenAI2026AreaLaw}.
\end{theorem}
\begin{proof}
    \leanok
    \uses{thm:replica_metric_comparison,thm:schur_label_commutation,
        lem:orthogonal_resolution,lem:orthogonal_resolution_cfc}
    Since $U\subseteq U\cup Y$, $Y\subseteq U\cup Y$, and
    $U\cap Y=\varnothing$, the three families of label projections
    commute. Their products form a common orthogonal resolution.
    On a nonzero product with labels $\alpha,\beta,\gamma$ for
    $U,U\cup Y,Y$, respectively, the squared metric factor has coefficient
    $q=(w_k(\gamma)/(w_k(\alpha)w_k(\beta)))^2$.
    Let $C_0\ge0$ be the constant in the logarithmic weight comparison
    of Theorem~\ref{thm:replica_metric_comparison}. Applying its lower
    bound to the first two weights and its upper bound to the third gives
    \begin{align*}
        \log q
         & \le 6C_0\log(k+2)
        +2t(\log d_\alpha+\log d_\beta-\log d_\gamma).
    \end{align*}
    All weights are positive. Exponentiating and summing over the common
    orthogonal resolution proves the assertion with $C=6C_0$.
    Zero products of label projections contribute nothing to either side.
\end{proof}

% Source: openai/math@adc7f1241b42e322a6451854ab7e4b4c146bf78a,
% 07-comparators.tex, lines 383--418. This is a leaf estimate only.

\begin{theorem}[Three signed regional moments]
    \label{thm:pinned_three_moment_holder}
    \lean{TensorPower.re_trace_mul_exp_pinnedLabels_le}
    \leanok
    Fix a finite family of finite-dimensional spaces and disjoint regions
    $U,Y$. On $k$ copies, write $F_Q$ for the central label entropy of
    region $Q$. For every positive semidefinite matrix $\rho$ on the whole
    copy space and every real $a$,
    \begin{align}
        \operatorname{Re}\operatorname{Tr}
            \bigl(\rho e^{a(F_U+F_{U\cup Y}-F_Y)}\bigr)
        &\leq
        \bigl(\operatorname{Re}\operatorname{Tr}(\rho e^{3aF_U})\bigr)^{1/3}
        \bigl(\operatorname{Re}\operatorname{Tr}(\rho e^{3aF_{U\cup Y}})\bigr)^{1/3}
        \bigl(\operatorname{Re}\operatorname{Tr}(\rho e^{-3aF_Y})\bigr)^{1/3}.
    \end{align}
    The matrix $\rho$ need not have trace one or commute with the three
    label observables. This is the three-factor inequality used in
    \cite[Section~8]{OpenAI2026AreaLaw}.
\end{theorem}
\begin{proof}
    \leanok
    \uses{thm:schur_label_commutation}
    The regions $U,Y,U\cup Y$ are pairwise nested or disjoint. Their
    central label observables therefore commute. Apply equal-weight trace
    H\"older to the three Hermitian matrices $F_U,F_{U\cup Y},-F_Y$.
\end{proof}

\begin{theorem}[Exterior exponential moments under Bell insertion]
    \label{thm:pinned_exponential_bell_insertion}
    \lean{TensorPower.re_inner_exp_labelEntropy_bellInsertion}
    \leanok
    Let $J_k$ insert the tensor power of a unit vector on the grouped
    physical region $X$ and auxiliary register $C$, with the original
    physical coordinates restored after insertion. The auxiliary register
    $R$ and the physical coordinates outside $X$ remain separate.
    Write $F_Q^{\mathrm{ext}}$ for the central label entropy on the
    exterior part of $Q$, and $F_{Q\setminus X}^{\mathrm{full}}$ for the
    corresponding operator in the original augmented space.
    For every exterior vector $\Psi$, real $u$, and physical region $Q$,
    \begin{align}
        \langle J_k\Psi,e^{uF_{Q\setminus X}^{\mathrm{full}}}J_k\Psi\rangle
        =\langle\Psi,e^{uF_Q^{\mathrm{ext}}}\Psi\rangle.
    \end{align}
\end{theorem}
\begin{proof}
    \leanok
    The central function $\lambda\mapsto e^{u\log\dim\lambda}$ gives
    $e^{uF_{Q\setminus X}^{\mathrm{full}}}J_k
      =J_ke^{uF_Q^{\mathrm{ext}}}$. Multiplication on the left by
    $J_k^*$ and $J_k^*J_k=I$ prove the identity. No metric weight is
    changed by this calculation.
\end{proof}

\begin{theorem}[A sharp upper bound for a pinned leaf]
    \label{thm:replica_pinned_sharp_leaf}
    \lean{TensorPower.exists_log_re_inner_replicaPinnedFactor_post_le}
    \leanok
    Fix a unit physical vector $\Omega$, a physical region $X$, and a
    spectral selection $E$ of its $X$ marginal with mass $z>0$.
    Let $\chi_E$ be its normalized compressed selected vector, with the
    selected register enumerated by $\{0,\ldots,|E|-1\}$, and let $b_E$
    be the normalized selected Bell vector on $X\otimes\mathbb C^E$.
    Set $p_k=J_k\chi_E^{\otimes k}$, where $J_k$ inserts $b_E^{\otimes k}$
    in the original physical coordinates. The original augmented family
    is the physical family together with $\mathbb C^E$ and
    $\mathbb C^{|E|}$.

    Fix $a>0$ and real $K,r$ with $6a\leq r$ and $6a\leq1$.
    There is a constant $C\geq0$, independent of $k,U,Y$, with the
    following property. Let $U,Y$ be disjoint physical regions, both
    disjoint from $X$. For $Q\in\{U,U\cup Y,Y\}$, let $\rho_Q$ be the
    original marginal of $\Omega$, and put $S_Q=S(\rho_Q)$. Suppose
    each marginal satisfies
    $\log M_{\rho_Q}(v)\leq vS_Q+Kv^2$ for all $|v|\leq r$,
    with zero spectral weights contributing zero to the moment.
    Write $W_Q$ for the replica metric of the original augmented family
    at parameter $a/2$, acting on the original physical copy of $Q$, and
    put $D_{U,Y}=(W_U^{-1}W_{U\cup Y}^{-1}W_Y)^2$.
    Then
    \begin{align}
        \log\langle p_k,D_{U,Y}p_k\rangle
        \leq{}&ka(S_U+S_{U\cup Y}-S_Y)+12Kka^2-k\log z\notag\\
        &+C\log(k+2)+\frac{(\dim H_Y)^2}{6}\log(k+1).
        \label{eq:pinned_sharp_leaf_rate}
    \end{align}
    The fixed constant $C$ may depend on the full augmented family and
    on $a$. Singular marginals, empty regions and zero copies are included.
    This is the leaf estimate in the sharp-upper proof of
    \cite[Section~8]{OpenAI2026AreaLaw}.
\end{theorem}
\begin{proof}
    \leanok
    \uses{thm:pinned_three_moment_holder,
        thm:pinned_exponential_bell_insertion,
        thm:replica_pinned_factor_exp_upper,
        thm:compressed_typical_exponential_domination,
        thm:iid_schur_exponential_moments}
    The original-family metric comparison, evaluated against
    $|p_k\rangle\langle p_k|$, bounds the expectation by
    $(k+2)^C\langle p_k,e^{a(F_U+F_{U\cup Y}-F_Y)}p_k\rangle$.
    Apply three-factor H\"older. The exterior insertion identity and
    spectral-selection domination bound each factor's moment by
    $z^{-k}$ times its original iid physical moment.

    The two positive moments at $3a$ have logarithmic bounds
    $3akS_U+9Kka^2$ and $3akS_{U\cup Y}+9Kka^2$.
    The negative moment has logarithmic bound
    $-3akS_Y+18Kka^2+(\dim H_Y)^2\log(k+1)/2$.
    Taking the three cube roots therefore gives quadratic coefficient
    $(9+9+18)/3=12$ and mass exponent $(-k-k-k)/3=-k$.
    Adding $C\log(k+2)$ proves (\ref{eq:pinned_sharp_leaf_rate}).
    The selected compressed vector is nonzero, the Bell insertion is
    isometric, and $D_{U,Y}$ is positive definite, so its expectation is
    strictly positive and taking its logarithm is legitimate.
\end{proof}

\begin{theorem}[The regions of successive bands]
\label{thm:replica_band_regions_subset_or_disjoint}
\lean{TensorPower.replicaBandRegions_subset_or_disjoint}
    \leanok
Let $F$ be a set and let $U_g,Y_g$ be finite subsets indexed by
$0\le g<K$. Suppose
\begin{align*}
 U_g\cap Y_g&=\varnothing, & U_g\cup Y_g&\subseteq U_h\quad(g<h).
\end{align*}
If $R$ is one of $U_g,Y_g,U_g\cup Y_g$ and $S$ is one of
$U_h,Y_h,U_h\cup Y_h$, then $R\subseteq S$, $S\subseteq R$, or
$R\cap S=\varnothing$.
\end{theorem}
\begin{proof}
    \leanok
For $g<h$, the set $R$ lies in $U_h$, so it is contained in $U_h$
and $U_h\cup Y_h$ and disjoint from $Y_h$. The case $h<g$ follows
by symmetry. For $g=h$, use the disjointness of $U_g,Y_g$ and their
inclusions in their union.
\end{proof}

\begin{theorem}[Replica metrics on a label cutoff]
    \label{thm:replica_metric_cutoff_norm}
    \lean{TensorPower.exists_replicaMetric_cutoff_norm_bound}
    \leanok
    \uses{def:schur_label_cutoff}
    Let the original system be a finite tensor product of nonzero
    finite-dimensional spaces, and fix $t>0$. Write $W_{Q,k}$ for its
    replica metric, using the same label weights on every subsystem.
    For $T\in\R$, let
    \begin{align*}
        C_{Q,k}(T)
         & =\sum_{\log d_\lambda\le T}\pi_{Q,k}^\lambda.
    \end{align*}
    There is $C\ge0$, depending on the original system and $t$, such
    that for every subsystem $Q$, every integer $k\ge0$, and every
    real threshold $T$,
    \begin{align*}
        \|W_{Q,k}\| & \le(k+2)^C,        \\
        \|W_{Q,k}^{-1}C_{Q,k}(T)\|
                    & \le(k+2)^C e^{tT}.
    \end{align*}
    These estimates supply the individual metric factors in the rough
    upper comparison of \cite[Section~8, upper comparisons]{OpenAI2026AreaLaw}. They retain the label weights of
    the original full system.
\end{theorem}
\begin{proof}
    \leanok
    \uses{thm:replica_metric_comparison,lem:orthogonal_resolution_cfc}
    On each nonzero label subspace, the replica-weight comparison gives
    \begin{align*}
        |\log w_k(\lambda)+t\log d_\lambda|
         & \le C\log(k+2).
    \end{align*}
    Since $d_\lambda\ge1$, this implies $w_k(\lambda)\le(k+2)^C$.
    If $\log d_\lambda\le T$, it also gives
    \begin{align*}
        w_k(\lambda)^{-1}
         & \le(k+2)^C e^{tT}.
    \end{align*}
    The metric, its inverse, and the cutoff are diagonal with respect
    to the same orthogonal label decomposition. Their operator norms
    are therefore bounded by these scalar coefficients. Labels absent
    from the representation contribute zero; labels outside the cutoff
    also contribute zero to the second operator.
\end{proof}

\begin{theorem}[Expectation of a pinned product leaf]
    \label{thm:replica_pinned_leaf_expectation}
    \lean{TensorPower.exists_replicaPinnedLeaf_expectation_le}
    \leanok
    Fix a finite tensor product of nonzero finite-dimensional spaces and
    $t>0$. Use the replica weights of this original full system to define
    $W_{Q,k}$ on every subsystem. There is $C\ge0$ such that the following
    holds for every $k,K\ge0$ and every real number $T$.

    Let $U_g,Y_g$, $0\le g<K$, be regions satisfying
    \begin{align*}
        U_g\cap Y_g&=\varnothing, &
        U_g\cup Y_g&\subseteq U_h\qquad(g<h).
    \end{align*}
    Suppose that a vector $\varphi$ lies in the closed label cutoffs
    $F_{U_g}\le T$ and $F_{U_g\cup Y_g}\le T$ for every $g$.
    Set
    \begin{align*}
        A&=\prod_{g=0}^{K-1}
          \bigl(W_{U_g,k}^{-1}W_{U_g\cup Y_g,k}^{-1}W_{Y_g,k}\bigr)^2,
    \end{align*}
    with the product taken in increasing order. Then
    \begin{align*}
        \operatorname{Re}\langle\varphi,A\varphi\rangle
        &\le (k+2)^{CK}\exp(4tKT)\|\varphi\|^2.
    \end{align*}
    The vector $\varphi$ is arbitrary in the stated cutoffs. This is the physical-region
    product in \cite[\texttt{07-comparators.tex},
    \texttt{comparator:pinned-metric}, lines 356--373]{OpenAI2026AreaLaw};
    its identification with the Bell restriction is separate.
\end{theorem}
\begin{proof}
    \leanok
    \uses{thm:replica_metric_cutoff_norm,thm:laminar_subsystem_label_commutation,thm:replica_band_regions_subset_or_disjoint}
    Write $C_0$ for the exponent in the metric cutoff estimate and put
    $p=(k+2)^{C_0}$ and $e=\exp(tT)$. The region relations make every
    pair among the regions $U_g,Y_g,U_g\cup Y_g$ nested or disjoint.
    Their central label functions commute. Hence each metric and inverse
    metric preserves the common cutoff space.

    On this space, the two inverse factors have norm at most $pe$,
    and the third factor has norm at most $p$. Each squared factor
    therefore has norm at most $p^6e^4$, and their ordered product has
    norm at most $p^{6K}e^{4K}$. Cauchy--Schwarz gives the stated
    expectation bound with $C=6C_0$. The same argument applies to the
    empty product and to zero copies. It also applies to negative $T$,
    for which the cutoff may have zero range.
\end{proof}


\section{Upper conditioning in a supplied sector}
\label{sec:mean_tree_supplied_projection_upper}

% Source: openai/math@adc7f1241b42e322a6451854ab7e4b4c146bf78a,
% 07-comparators.tex, lines 383--404.

\begin{theorem}[Quadratic contraction by a commuting projection]
    \label{thm:commuting_projection_quadratic_upper}
    \lean{Matrix.PosSemidef.re_dotProduct_projection_mulVec_le}
    \leanok
    Let $A$ be positive semidefinite and $P$ an orthogonal projection
    commuting with $A$. For every vector $p$,
    \[
        \langle Pp,APp\rangle\leq\langle p,Ap\rangle.
    \]
\end{theorem}
\begin{proof}
    \leanok
    The complementary sandwich satisfies
    $(I-P)^*A(I-P)=A-PAP\geq0$. Evaluate this inequality at $p$,
    using $P^*=P$ to identify the compressed quadratic form.
\end{proof}

\begin{theorem}[A terminal-weighted upper bound in a supplied sector]
    \label{thm:mean_tree_supplied_projection_upper}
    \lean{Matrix.MeanTree.re_inner_log_eval_le_of_projection}
    \leanok
    Let $T$ be a finite weighted mean tree with positive definite leaf
    matrices $A_j$ and terminal weights $w_j$. Let $P$ be a supplied
    orthogonal projection commuting with every $A_j$, and let $p$ be
    a supplied vector with $q=\|Pp\|^2>0$. Put $\phi=Pp/\|Pp\|$.
    Then
    \[
        \langle\phi,\log(T(A))\phi\rangle
        \leq-\log q+\sum_jw_j\log\langle p,A_jp\rangle.
    \]
    The original vector need not have norm one. The leaves need not
    commute with one another. Repeated labels and zero terminal weights
    remain included.
\end{theorem}
\begin{proof}
    \leanok
    \uses{thm:commuting_projection_quadratic_upper}
    The projection commutes with the root $M=T(A)$, which is positive
    definite. Vector log Jensen, quadratic contraction at that root,
    and vector mean Jensen give
    \[
        \langle\phi,\log M\,\phi\rangle
        \leq\log\langle\phi,M\phi\rangle
        \leq\log\left(q^{-1}\langle p,Mp\rangle\right)
        \leq\log\left(q^{-1}
            \prod_j\langle p,A_jp\rangle^{w_j}\right).
    \]
    Every logarithm has a positive argument, since $q>0$ implies $p\ne0$.
    Expanding the scalar product logarithm gives the asserted bound.
    This retains the exact conditioning cost in
    \cite[Section~8]{OpenAI2026AreaLaw}.
\end{proof}

\begin{theorem}[The supplied mass in a one-band inverse norm]
    \label{thm:mean_tree_projection_inverse_norm_upper_consumer}
    \lean{Matrix.MeanTree.mass_rpow_mul_exp_le_norm_rpow_neg_eval_projection_sq}
    \leanok
    Under the preceding hypotheses, for every real $s\geq0$,
    \[
        q^{1+2s}\exp\left(-2s\sum_jw_j
            \log\langle p,A_jp\rangle\right)
        \leq\|T(A)^{-s}Pp\|^2.
    \]
    This is a one-band statement. It neither selects the projection
    nor supplies a lower bound on its mass.
\end{theorem}
\begin{proof}
    \leanok
    \uses{thm:mean_tree_supplied_projection_upper}
    Apply inverse-power exponential Jensen to the literal normalized
    vector $Pp/\|Pp\|$ and then undo normalization. The retained factor
    $q$ and the conditioning term $-\log q$ give
    $q\exp(2s\log q)=q^{1+2s}$.
\end{proof}

\section{Sharp bounds in the same supplied sector}
\label{sec:selected_bell_post_sharp_tree}

% Static source candidate; no completion marks before successful verification.
% Source: openai/math@adc7f1241b42e322a6451854ab7e4b4c146bf78a,
% 07-comparators.tex, lines 383--418.

\begin{theorem}[One supplied sector across commuting bands]
    \label{thm:mean_tree_projection_multiband_log_upper}
    \lean{Matrix.MeanTree.re_inner_log_eval_prod_le_of_projection}
    \leanok
    For each $g\in\{1,\ldots,G\}$, let $T_g$ be a finite weighted mean
    tree with positive definite leaves $A_{j,g}$ and terminal weights
    $w_{j,g}$. Suppose its roots $M_g=T_g(A_{\cdot,g})$ commute across
    distinct bands. Let $P$ be a supplied orthogonal projection commuting
    with every leaf, and let $p$ satisfy $q=\|Pp\|^2>0$.
    Put $M=\prod_{g=1}^G M_g$ in band order and $\phi=Pp/\|Pp\|$.
    Then
    \begin{align}
        \langle\phi,\log M\,\phi\rangle
        &\leq-G\log q+
            \sum_{g=1}^G\sum_jw_{j,g}\log\langle p,A_{j,g}p\rangle.
    \end{align}
    No normalization of $p$ or mutual commutation of leaves in a single
    band is required. Empty band products are included.
\end{theorem}
\begin{proof}
    \leanok
    \uses{thm:mean_tree_supplied_projection_upper}
    Positivity and commutation of the roots give
    $\log M=\sum_g\log M_g$. Apply the supplied-sector logarithmic
    bound to each root, always with the same $P$, $p$ and $q$, and sum.
    Each band contributes exactly $-\log q$.
\end{proof}

\begin{theorem}[The supplied mass in a multi-band inverse norm]
    \label{thm:mean_tree_projection_multiband_inverse_upper}
    \lean{Matrix.MeanTree.mass_rpow_mul_exp_le_norm_rpow_neg_eval_prod_projection_sq}
    \leanok
    Under the hypotheses of
    Theorem~\ref{thm:mean_tree_projection_multiband_log_upper}, every real
    $s\geq0$ satisfies
    \begin{align}
        q^{1+2sG}\exp\left(-2s\sum_{g=1}^G\sum_jw_{j,g}
            \log\langle p,A_{j,g}p\rangle\right)
        &\leq\|M^{-s}Pp\|^2.
    \end{align}
\end{theorem}
\begin{proof}
    \leanok
    \uses{thm:mean_tree_projection_multiband_log_upper}
    Exponential Jensen in $Pp/\|Pp\|$ gives
    $q\exp(-2s\langle\phi,\log M\,\phi\rangle)
        \leq\|M^{-s}Pp\|^2$.
    Substitute the logarithmic bound and combine
    $q\exp(2sG\log q)=q^{1+2sG}$.
\end{proof}

\begin{definition}[The selected Bell postvector]
    \label{def:selected_bell_postvector}
    \lean{TensorPower.selectedBellPost}
    \leanok
    Let $\Omega$ be a physical vector, let $X$ be a region, and let $E$
    be a finite spectral selection of its reduced state on $X$.
    Write $\Omega_X$ in the coordinates $H_X\otimes H_{X^c}$, let
    $\chi_E$ be its selected compression with inverse-square-root mass
    normalization and the selected
    eigenvectors enumerated by $\{1,\ldots,|E|\}$, and let $b_E$ be
    the corresponding selected Bell vector. In the original physical
    coordinates augmented by $E$ and $\{1,\ldots,|E|\}$, define
    $\mathrm{post}_k=J_{b_E}\chi_E^{\otimes k}$, where $J_{b_E}$ is
    the literal Bell insertion, with the physical coordinates unchanged.
    The algebraic definitions also apply at zero selected mass and zero copies.
\end{definition}

\begin{theorem}[Terminal-weighted sharp bounds for the selected postvector]
    \label{thm:selected_bell_post_sharp_tree}
    \lean{TensorPower.exists_selectedBellPost_sharp_tree_bounds}
    \leanok
    Fix a finite physical family, a unit vector $\Omega$, a region $X$,
    a spectral selection $E$ of positive mass $z$, and real parameters
    $a>0$, $K$, $r$ with $6a\leq r$ and $6a\leq1$.
    There is a nonnegative constant $C$, chosen before the copy number,
    bands, terminal labels, regions and weighted tree, with the following property.

    Let $U_{j,g},Y_{j,g}$ be disjoint from $X$ and each other, with
    $U_{j,g}\cup Y_{j,g}\subseteq U_{j',h}$ whenever $g<h$.
    Suppose each original physical marginal on
    $Q\in\{U_{j,g},U_{j,g}\cup Y_{j,g},Y_{j,g}\}$ obeys the one-copy
    surprisal bound
    \begin{align}
        \log\operatorname{Tr}\rho_Q^{1-v}
        &\leq vS_\Omega(Q)+Kv^2
    \end{align}
    for every $|v|\leq r$, with the support convention for singular marginals.
    Let $T$ be a finite weighted mean tree with terminal weights $w_j$,
    let $W$ use the original full-family replica weights at $a/2$, and
    write $\operatorname{ext}Q$ for the physical coordinates of $Q\setminus X$
    in the augmented family. Define the actual exterior factors by
    \begin{align}
        D_{j,g}
        &=\left(W_{\operatorname{ext}U_{j,g}}^{-1}
            W_{\operatorname{ext}(U_{j,g}\cup Y_{j,g})}^{-1}
            W_{\operatorname{ext}Y_{j,g}}\right)^2.
    \end{align}
    Put
    $M=\prod_gT(D_{\cdot,g})$ in band order.
    For the supplied auxiliary label $\lambda$, set
    $P=\pi_R^\lambda$, $q=\|P\mathrm{post}_k\|^2$ and suppose $q>0$.
    Define
    \begin{align}
        L_k
        &=ka\sum_g\sum_jw_j
            \bigl(S_\Omega(U_{j,g})+S_\Omega(U_{j,g}\cup Y_{j,g})
                -S_\Omega(Y_{j,g})\bigr)\notag\\
        &\quad+G\bigl(12Kka^2-k\log z+C\log(k+2)\bigr).
    \end{align}
    For $\phi=P\mathrm{post}_k/\|P\mathrm{post}_k\|$ and every real $s\geq0$,
    \begin{align}
        \langle\phi,\log M\,\phi\rangle&\leq-G\log q+L_k,\\
        q^{1+2sG}\exp(-2sL_k)&\leq\|M^{-s}P\mathrm{post}_k\|^2.
    \end{align}
    These finite-copy bounds retain the supplied label and its exact mass.
    They do not select the label, give a polynomial lower bound for $q$,
    or identify the inverse norm with the original physical comparison vector.
\end{theorem}
\begin{proof}
    \leanok
    \uses{def:selected_bell_postvector,thm:replica_pinned_sharp_leaf,
        thm:mean_tree_projection_multiband_log_upper,
        thm:mean_tree_projection_multiband_inverse_upper}
    Unit norm and positive selected mass give the nonempty coordinate
    factors required for positivity. Nesting gives positive definite
    pinned factors and cross-band commutation. A one-site central label
    projection commutes with every subsystem label observable, hence with
    the actual metric factors and their inverses.

    Let $C_0$ be the uniform replica-weight coefficient in the sharp
    leaf bound and choose
    $C=C_0+\sum_{Q\subseteq V}(\dim H_Q)^2/6$.
    Since $\log(k+1)\leq\log(k+2)$, this one fixed coefficient absorbs
    every leaf's dimensional remainder. Apply the sharp leaf estimate
    to the literal same $\Omega$ and $E$, multiply by $w_j\geq0$ and
    sum. The identity $\sum_jw_j=1$ leaves one cost $-k\log z$ per band.
    The two multi-band inequalities then give the conclusions without
    changing the projection, its mass, or any terminal weight.
\end{proof}

\section{Exact mass and physical transport in the selected sector}
\label{sec:selected_bell_sector_transport}

\begin{theorem}[The same auxiliary label through Bell insertion]
    \label{thm:selected_bell_auxiliary_intertwining}
    \lean{TensorPower.labelProj_mul_bellInsertion_auxiliary,TensorPower.norm_labelProj_bellInsertion_auxiliary}
    Let $J_b$ insert $b^{\otimes k}$ on the physical region $X$ and
    auxiliary register $C$, retaining the distinct auxiliary register $R$.
    Let $Q_\lambda$ and $Q^e_\lambda$ be the central projections of the
    same irreducible label $\lambda$ on the original and exterior $R$
    copies. Then $Q_\lambda J_b=J_bQ^e_\lambda$.
    If $\|b\|=1$, every exterior vector $\psi$ satisfies
    $\|Q_\lambda J_b\psi\|=\|Q^e_\lambda\psi\|$.
    These assertions include zero copies and empty coordinate spaces.
\end{theorem}
\begin{proof}
    \uses{thm:bell_insertion_label_intertwining,thm:bell_insertion_isometry_symmetry}
    Apply the auxiliary-label observable identity to the indicator
    function $\mu\mapsto\mathbf1_{\{\mu=\lambda\}}$.
    A unit inserted vector gives $J_b^*J_b=I$, so
    $\|Q_\lambda J_b\psi\|=\|J_bQ^e_\lambda\psi\|
      =\|Q^e_\lambda\psi\|$.
\end{proof}

Fix a physical vector $\Omega$, a region $X$, and a set $E$ of spectral
indices for its actual marginal on $X$. Put $d=|E|$, and retain the fixed
enumeration of $E$ by $\{0,\ldots,d-1\}$. The left auxiliary register is
$C=E$; the right register is $R=\{0,\ldots,d-1\}$.
Let $B$ be the physical coordinate set on $X^c$.
Write $b_E$ for the selected Bell vector, $\chi_E$ for the selected
compressed vector, and $\Psi_k$ for its tensor power in the
exterior sites $R\cup X^c$. Thus $\mathrm{post}_k=J_{b_E}\Psi_k$ is the
vector of Definition~\ref{def:selected_bell_postvector}.

\begin{theorem}[Exact mass in the supplied auxiliary sector]
    \label{thm:selected_bell_post_mass}
    \lean{TensorPower.norm_sq_labelProj_selectedBellPost_eq,TensorPower.norm_sq_labelProj_selectedBellPost_prod}
    Suppose the actual selected mass $z$ is positive. Let
    $\psi(b,r)=\chi_E(\iota^{-1}(r),b)$ put the physical complement first,
    and group its $k$ physical factors before its $k$ auxiliary factors.
    For every $k\ge0$ and every supplied label $\lambda$,
    \begin{align*}
        \|Q_\lambda\mathrm{post}_k\|^2
        &=\|Q^e_\lambda\Psi_k\|^2
          =\|(I_{B^k}\otimes\Pi^{d,k}_\lambda)\psi^{\otimes k}\|^2.
    \end{align*}
    No normalization of $\Omega$ is required. In particular, the equality
    holds pointwise along any already chosen label sequence.
\end{theorem}
\begin{proof}
    \uses{thm:selected_bell_auxiliary_intertwining,thm:compressed_site_label_mass_coordinates}
    Positive selected mass makes $b_E$ unit, and the insertion preserves
    the norm of its supplied auxiliary sector. Regrouping the exterior
    coordinates is unitary and carries the same auxiliary projection to
    $I_{B^k}\otimes\Pi^{d,k}_\lambda$. Neither operation changes the label.
\end{proof}

\begin{theorem}[The postvector in separated Bell coordinates]
    \label{thm:selected_bell_post_pullback}
    \lean{TensorPower.selectedBellPost_eq_pullback}
    Let $\tau_k$ separate the selected Bell coordinates from the
    exterior coordinates, and let $W_kf=f\circ\tau_k$.
    The literal postvector satisfies
    \begin{align*}
        \mathrm{post}_k
        &=W_k\bigl(b_E^{\otimes k}\otimes\Psi_k\bigr).
    \end{align*}
    This coordinate equality requires no normalization or positive-mass
    hypothesis and includes zero copies.
\end{theorem}
\begin{proof}
    \uses{def:selected_bell_postvector,def:schmidt_bell_witness_coordinates,
        thm:added_site_tensor_insertion_action}
    Applying the insertion matrix multiplies the exterior coefficient by
    $\prod_jb_E(x_j|_X,c_j)$. The remaining coefficient is precisely the
    product defining $\Psi_k$ in the exterior coordinates of $\tau_k$.
\end{proof}

\begin{theorem}[Symmetry of the supplied postvector sector]
    \label{thm:selected_bell_post_sector_symmetry}
    \lean{TensorPower.symProj_labelProj_selectedBellPost}
    Let $S_k$ average simultaneous permutations of the original copies.
    For every supplied label $\lambda$,
    $S_kQ_\lambda\mathrm{post}_k=Q_\lambda\mathrm{post}_k$.
    No normalization or positive selected-mass hypothesis is required.
\end{theorem}
\begin{proof}
    \uses{thm:selected_bell_auxiliary_intertwining,
        thm:bell_insertion_isometry_symmetry,thm:central_label_full_symmetry}
    If $S^e_k$ is exterior copy symmetry, the pure tensor satisfies
    $S^e_k\Psi_k=\Psi_k$ and the central auxiliary projection commutes
    with $S^e_k$. Since $S_kJ_{b_E}=J_{b_E}S^e_k$,
    \begin{align*}
        S_kQ_\lambda J_{b_E}\Psi_k
        &=J_{b_E}S^e_kQ^e_\lambda\Psi_k
          =J_{b_E}Q^e_\lambda\Psi_k
          =Q_\lambda J_{b_E}\Psi_k.
    \end{align*}
\end{proof}

\begin{theorem}[The exact symmetric Bell pin]
    \label{thm:selected_bell_post_exact_pin}
    \lean{TensorPower.bellPin_schmidtBellPrevector_eq_selectedBellPost}
    Suppose $z>0$. Let $v_{k,\lambda}$ be the original physical prevector
    with its supplied auxiliary label, expressed in the original physical
    sites, selected left register $E$, and enumerated right register $R$.
    Set $\mathsf P_k=J_{b_E}S^e_kJ_{b_E}^*$. Then
    \begin{align*}
        \mathsf P_kv_{k,\lambda}
        &=\left(\frac{\sqrt z}{d}\right)^k
            Q_\lambda\mathrm{post}_k.
    \end{align*}
    The state, selected set, enumeration and label are the same on both
    sides. This holds for every $k\ge0$, without a unit-norm condition on
    $\Omega$, and is the symmetric-pin form of
    \cite[\texttt{07-comparators.tex}, lines 283--290]{OpenAI2026AreaLaw}.
\end{theorem}
\begin{proof}
    \uses{thm:schmidt_bell_site_pin,thm:schmidt_bell_witness_outer_product,
        thm:schmidt_bell_witness_auxiliary_label,thm:selected_bell_post_pullback,
        thm:selected_bell_post_sector_symmetry,thm:bell_insertion_isometry_symmetry}
    Pull back the selected Bell identity by $\tau_k$. The two projection
    identities and the literal postvector equality give
    $J_{b_E}J_{b_E}^*v_{k,\lambda}
      =(\sqrt z/d)^kQ_\lambda\mathrm{post}_k$.
    Insertion symmetry gives
    $\mathsf P_k=S_kJ_{b_E}J_{b_E}^*$.
    The supplied postvector sector is fixed by $S_k$, proving the claim.
\end{proof}

\begin{theorem}[Physical inverse norm with the exact pin cost]
    \label{thm:selected_bell_post_inverse_transport}
    \lean{TensorPower.selectedBellPost_inverse_norm_le_physical}
    Assume $\|\Omega\|=1$ and $z>0$. Fix $t\ge0$, $k\ge0$, and an
    isometry $Z:\C^n\to\mathcal H_k$ with $ZZ^*=S_k$.
    Let $U_{jg},Y_{jg}$ satisfy $U_{jg}\cap Y_{jg}=\varnothing$ and,
    for every $g<h$, $U_{jg}\cup Y_{jg}\subseteq U_{j'h}$.
    Define $M$ and $N$ as in
    Theorem~\ref{thm:bell_compressed_band_product_powers}, using the same
    weighted mean tree in every band and the original full-family
    replica weights. For every supplied label $\lambda$ and every
    real $s$,
    \begin{align*}
        \left(\frac z{d^2}\right)^k
          \|N^{-s}Q_\lambda\mathrm{post}_k\|^2
        &\le\|M^{-s}Z^*v_{k,\lambda}\|^2.
    \end{align*}
    Empty band products and zero copies are included. This is the
    finite-copy pin-norm comparison used in
    \cite[\texttt{07-comparators.tex}, lines 283--308 and 383--418]{OpenAI2026AreaLaw}.
\end{theorem}
\begin{proof}
    \uses{thm:selected_bell_post_exact_pin,thm:bell_compressed_band_product_powers,
        thm:bell_band_product_pin_action,thm:bell_symmetric_pin_projection,
        thm:projection_intertwining_norm}
    The unit physical state and positive selected mass imply nonempty
    original coordinate sets. The actual band-product theorems give
    positive $M,N$ and
    $(ZM^{-s}Z^*)\mathsf P_k=\mathsf P_kN^{-s}$.
    Since $b_E$ is unit, $\mathsf P_k$ is an orthogonal projection.
    The projected norm comparison and the exact pin identity yield
    \begin{align*}
        \left|\left(\frac{\sqrt z}{d}\right)^k\right|
          \|N^{-s}Q_\lambda\mathrm{post}_k\|
        &=\|N^{-s}\mathsf P_kv_{k,\lambda}\|
          \le\|M^{-s}Z^*v_{k,\lambda}\|.
    \end{align*}
    Squaring and using $|\sqrt z/d|^{2k}=(z/d^2)^k$ proves the result.
\end{proof}

\section{Uniform quadratic-form bounds through a mean tree}
\label{sec:mean_tree_uniform_expectation}

Let $\mathcal T$ be a finite binary tree of weighted geometric means, with
positive definite leaf operators on a finite-dimensional complex Hilbert space $(A_j)_{j\in J}$ indexed by a finite set $J$,
and root $M$.
Write $w_j\ge0$ for its terminal weights, so $\sum_jw_j=1$.
The following consequences of vector Jensen and spectral Jensen are used
in \cite[Section~8, upper comparisons]{OpenAI2026AreaLaw}.

\begin{theorem}[A uniform leaf bound]
    \label{thm:mean_tree_uniform_expectation}
    \lean{Matrix.MeanTree.re_inner_eval_le_of_forall}
    \leanok

    Let $L>0$ and let $\varphi$ be any vector. If
    $\langle\varphi,A_j\varphi\rangle\le L$ for every $j$, then
    $\langle\varphi,M\varphi\rangle\le L$.
\end{theorem}
\begin{proof}
    \leanok
    \uses{thm:mean_tree_properties}
    Vector Jensen and monotonicity of nonnegative real powers give
    \begin{align*}
        \langle\varphi,M\varphi\rangle
        &\le\prod_j\langle\varphi,A_j\varphi\rangle^{w_j}
        \le\prod_j L^{w_j}=L.
    \end{align*}
\end{proof}

\begin{theorem}[A uniform bound after a fractional power]
    \label{thm:mean_tree_uniform_power_expectation}
    \lean{Matrix.MeanTree.re_inner_rpow_eval_le_of_forall}
    \leanok
    Under the same hypotheses, suppose additionally that
    $\|\varphi\|=1$. For every $p\in[0,1]$,
    $\langle\varphi,M^p\varphi\rangle\le L^p$.
\end{theorem}
\begin{proof}
    \leanok
    \uses{thm:mean_tree_uniform_expectation,lem:weighted_geometric_mean_vertex}
    Apply concavity of $x\mapsto x^p$ to the spectral probability measure
    of $M$ at $\varphi$, and then apply the uniform root bound:
    \begin{align*}
        \langle\varphi,M^p\varphi\rangle
        &\le\langle\varphi,M\varphi\rangle^p\le L^p.
    \end{align*}
\end{proof}

\begin{theorem}[An overlap bound for an inverse power]
    \label{thm:mean_tree_uniform_inverse_overlap}
    \lean{Matrix.MeanTree.norm_inner_sq_le_rpow_mul_inverse_norm_sq}
    \leanok
    Under the same hypotheses with $\|\varphi\|=1$, for every
    $s\in[0,1/2]$ and every vector $u$,
    \begin{align*}
        |\langle\varphi,u\rangle|^2
        &\le L^{2s}\|M^{-s}u\|^2.
    \end{align*}
\end{theorem}
\begin{proof}
    \leanok
    \uses{thm:mean_tree_uniform_power_expectation}
    Cauchy--Schwarz with opposite powers gives
    $|\langle\varphi,u\rangle|^2
      \le\langle\varphi,M^{2s}\varphi\rangle\|M^{-s}u\|^2$.
    Apply the uniform power bound with $p=2s$.
\end{proof}

\begin{theorem}[The actual product of band roots]

    \label{thm:replica_pinned_tree_band_product}
    \lean{TensorPower.replicaPinnedTree_eval_eq_prod}
    \leanok
    Fix a finite tensor product of nonzero finite-dimensional spaces and
    $t\ge0$. Use its original full-system weights to define the replica
    metrics $W_{Q,k}$. Let $\mathcal T$ be one weighted geometric-mean tree
    with leaf labels in $J$. Suppose the regions $U_{j,g},Y_{j,g}$ satisfy
    \begin{align*}
        U_{j,g}\cap Y_{j,g}&=\varnothing, &
        U_{j,g}\cup Y_{j,g}&\subseteq U_{j',h}
        \qquad(g<h)
    \end{align*}
    for every $j,j'$. Set
    \begin{align*}
        D_{j,g}
        &=\bigl(W_{U_{j,g},k}^{-1}
          W_{U_{j,g}\cup Y_{j,g},k}^{-1}W_{Y_{j,g},k}\bigr)^2,
        & M_g&=\mathcal T(D_{\cdot,g}).
    \end{align*}
    Then, with all products in increasing band order,
    \begin{align*}
        \mathcal T\left(\left(\prod_{g<K}D_{j,g}\right)_{j\in J}\right)
        &=\prod_{g<K}M_g.
    \end{align*}
    This is the actual band factorization in
    \cite[\texttt{06-transport.tex}, Lemma~7.1;
    \texttt{07-comparators.tex}, lines 367--375]{OpenAI2026AreaLaw}.
\end{theorem}
\begin{proof}
    \leanok
    \uses{thm:laminar_replica_band_commutation,thm:mean_tree_properties}

    Each $D_{j,g}$ is positive definite by the one-band case of the
    preceding positivity result. If $g<h$, every region in band $g$
    lies inside $U_{j',h}$. It is consequently nested with $U_{j',h}$
    and $U_{j',h}\cup Y_{j',h}$ and disjoint from $Y_{j',h}$.
    The central label functions therefore commute across distinct bands,
    even for different leaf labels. Apply the factorization of
    a weighted mean tree across commuting bands.
\end{proof}

\begin{theorem}[An inverse-power estimate for the actual band product]

    \label{thm:replica_pinned_tree_inverse_overlap}
    \lean{TensorPower.exists_replicaPinnedTree_inverse_norm_bound}
    \leanok
    Fix the original finite tensor product, $t>0$, and a finite leaf-label
    set $J$. There is $C\ge0$ such that the following holds for every copy
    number $k$, every number $K$ of bands, and every weighted mean tree
    $\mathcal T$ on $J$.

    Let the regions satisfy the preceding disjointness and cross-label
    nesting conditions, and put $M=\prod_{g<K}M_g$. Let $a\in\R$
    and let $\varphi$ be a unit vector in all the closed cutoffs
    $F_{U_{j,g}}\le a$ and $F_{U_{j,g}\cup Y_{j,g}}\le a$.
    For every vector $u$ and every $0\le s\le1/2$,
    \begin{align*}
        |\langle\varphi,u\rangle|^2
        &\le
        \bigl((k+2)^{CK}\exp(4tKa)\bigr)^{2s}
        \|M^{-s}u\|^2.
    \end{align*}
    The constant is chosen before $k,K$, the regions, the tree, the
    cutoff, and the vectors. The statements include $k=0$, $K=0$,
    and negative cutoff thresholds. This is the inverse-power step of
    \cite[\texttt{07-comparators.tex}, lines 367--381]{OpenAI2026AreaLaw}.
\end{theorem}
\begin{proof}
    \leanok
    \uses{thm:replica_pinned_tree_band_product,thm:replica_pinned_leaf_expectation,thm:mean_tree_uniform_inverse_overlap,thm:positive_laminar_replica_product}

    The homogeneous leaf estimate gives the same positive bound
    $L=(k+2)^{CK}\exp(4tKa)$ for the quadratic form at every ordered
    leaf product. Each such product is positive definite. Vector Jensen,
    spectral Jensen, and Cauchy--Schwarz with opposite powers yield
    $|\langle\varphi,u\rangle|^2\le L^{2s}\|\mathcal T(A)^{-s}u\|^2$.
    The preceding exact factorization replaces this root by the actual
    ordered product $M$ of the band roots.
\end{proof}
